\begin{abstract}
Large vision--language models (LVLMs) frequently hallucinate objects that are absent from the visual input, limiting their reliability in practical applications. Training-free contrastive decoding can reduce such errors, but many existing methods require multiple model evaluations per decoding step. The method ONLY mitigates hallucinations within a single forward process by suppressing attention heads selected from a Text-to-Visual Entropy Ratio (TVER) at one intervention layer and propagating the resulting representation through a lightweight ghost path. However, a mask computed at one early layer cannot incorporate deeper cross-modal evidence, and a single masked branch provides only one distorted view of the model's internal behavior.

We introduce \emph{Adaptive Dual-Branch Decoding} (ADBD), a training-free extension of ONLY that accumulates continuous TVER evidence across selected Transformer layers using head-specific adaptive gains. Rather than committing to a binary decision at one layer, ADBD maintains an evolving state for each head and adjusts the update magnitude according to the confidence and innovation of newly observed evidence. The accumulated state induces two complementary masks, yielding high-TVER and low-TVER auxiliary branches alongside the normal branch. At each autoregressive step, ADBD independently measures the Total Variation Distance (TVD) between each auxiliary distribution and the normal distribution, then adaptively reinforces or penalizes each branch through an anchored three-branch decoding rule.
\end{abstract}
